\section{Descent before conversion}
\label{sec:descent}

We now combine the two principal theorems. The geometric input describes
a finite cover while retaining its residual torus action. Invariant forms
on a compact factor then descend to that cover, where the ordinary
conversion theorem applies.

We use the following finite-cover structure statement from
\cite[Theorem~(Finite cover with compact torus monodromy)]{CompanionH}. If \(X\) is smooth
connected projective and \(-K_X\) is smoothly semipositive, there is
a connected finite \'etale projective cover
\(\nu:S\to X\) whose universal cover is
\[
\widetilde S=\C^b\times C\times F.
\]
Here \(C\) and \(F\) are compact simply connected projective factors;
\(C\) is Ricci-flat, and \(F\) has smoothly semipositive anticanonical
bundle and no positive-degree holomorphic forms. The deck group acts by
translations on \(\C^b\), trivially on \(C\), and through a compact
real torus \(T\) of holomorphic isometries on \(F\). The first two
factors have nowhere-zero holomorphic canonical forms invariant under
the deck group. Empty products are points and have unit canonical frame.

The proof of this statement in the companion starts from Yau's
prescribed-Ricci theorem and the Ricci-semipositive structure theory
\cite{Yau1978,DPS1996,CDP2015}. Its finite-cover construction preserves
monodromy and the invariant frames. It also proves the vanishing of
holomorphic forms by the Bochner argument on irreducible
non-Ricci-flat factors, independently of the rational-connectedness
conclusion of the refined structure theorem. K\"ahler Hodge symmetry
therefore gives
\begin{equation}\label{eq:factor-euler}
H^q(F,\OO_F)=0\quad(q>0),\qquad H^0(F,\OO_F)=\C,
\qquad I_T(0)=1.
\end{equation}
This deduction includes a point factor.

\begin{proposition}[Descent of invariant twisted forms]\label{prop:descent}
Let \(\nu:S\to X\), \(F\), and \(T\) be the finite-cover data just
specified. After choosing invariant canonical frames on \(\C^b\) and
\(C\), there is, for every \(r\ge0\) and \(0\le p\le\dim F\), an
injective linear map
\[
H^0(F,\Omega_F^p\otimes(-K_F)^r)^T
\longrightarrow H^0(S,\Omega_S^p\otimes(-K_S)^r).
\]
For some fixed \(p\), the target is nonzero at unbounded positive
integers \(r\).
\end{proposition}

\begin{proof}
Let \(\eta\) be the product of the chosen canonical frames on
\(\C^b\) and \(C\). For an invariant form \(u\) in the displayed
source, pull back its differential part along the projection to \(F\)
and use \(\eta^{-r}\) for the other anticanonical factors. This gives
\[
\eta^{-r}\otimes\operatorname{pr}_F^*u
\in H^0\bigl(\widetilde S,
\Omega_{\widetilde S}^p\otimes K_{\widetilde S}^{-r}\bigr).
\]
Translations preserve the Euclidean frame, the action on \(C\) is
trivial, and the action on \(F\) lies in \(T\), with its natural induced
action on both tensor factors of \(u\). The tensor is therefore
deck-invariant and descends holomorphically to \(S\). Pullback along
the projection is injective on differential forms, tensoring with
\(\eta^{-r}\) is invertible, and local covering charts preserve
nonvanishing. Hence the descent map is injective. Its holomorphic
sections are algebraic because \(S\) is projective.

By \eqref{eq:factor-euler} and Corollary~\ref{cor:invariantforms}, there
are a fixed \(p\) and invariant sections in
\(H^0(F,\Omega_F^p\otimes(-K_F)^{m+1})\) for unbounded positive
\(m\). Apply the constructed map with \(r=m+1\). If \(F\) is a
point, its constant section gives the same assertion with \(p=0\).
\end{proof}

\begin{corollary}[Smooth anticanonical nonvanishing]\label{cor:global}
Let \(X\) be a smooth connected projective complex variety. If
\(-K_X\) has a smooth Hermitian metric with semipositive Chern
curvature, then \(H^0(X,-mK_X)\ne0\) for some integer \(m>0\).
\end{corollary}

\begin{proof}
Choose the finite cover above. Proposition~\ref{prop:descent} gives
nonzero anticanonically twisted differential forms on \(S\) in one
fixed degree and at unbounded positive exponents. The anticanonical
metric on \(S\) is the smooth semipositive pullback metric. Apply
Corollary~\ref{cor:conversion} on \(S\) to obtain a nonzero section of
\(-kK_S\), \(k>0\).

The finite \'etale norm of a section
\cite[Lemma~(Finite \'etale norm)]{CompanionH} sends a nonzero section of
\(\nu^*A\) to a nonzero section of \(A^{\otimes d}\), where
\(d=\deg\nu\); a Galois hypothesis is unnecessary. With
\(A=-kK_X\), the identity \(K_S=\nu^*K_X\) therefore gives a section
of \(-dkK_X\). For zero-dimensional connected \(X\), the canonical
line is trivial and the assertion is immediate.
\end{proof}

The ordering of the argument matters. The torus-invariant objects are
the forms on \(F\); they descend before any section-conversion theorem
is applied. Consequently no equivariant strengthening of
Theorem~\ref{thm:conversion} is required, and the finite cover is never
treated as a global product.
